% ===== 原文 pp.27–28 =====
% （§4.3.3 后半；图 11；本块 tag 4.22–4.27；收在 p.28 锚点句为止，§5 标题归 tr-05）

% [原文 p.27]
\begin{center}
\begin{minipage}{\textwidth}
\centering
\includegraphics[width=0.45\textwidth]{fig11.png}\par\vspace{2pt}
\figtitle{11}{Young 图 $Y=(5,4,3,2,2,1)$ 与盒子 $s=(2,2)$，以及它的臂长与leg长。}
\end{minipage}
\end{center}

现在我们可以最终写下瞬子贡献了。例如，考虑一个带有双基表示超多重态的 $\SU(N)\times\SU(M)$ 规范理论。它的瞬子贡献是对可能的点状瞬子组态求和：
\begin{equation}
Z_{\text{instanton}}(\vec a;\vec b)=\sum_{\vec Y,\vec W}q^{|\vec Y|}q'^{|\vec W|}Z(\vec a,\vec Y;\vec b,\vec W)
\tag{4.22}
\end{equation}
其中 $\vec a=(a_1,\ldots,a_N)$ 与 $\vec b=(b_1,\ldots,b_M)$ 是 $\SU(N)$ 与 $\SU(M)$ 矢量多重态中伴随标量实部的真空期望值，而 $\vec Y=(Y_1,\ldots,Y_N)$、$\vec W=(W_1,\ldots,W_M)$ 标记点状瞬子组态。前置因子 $q,q'$ 由 $q=\ee^{2\pi\sqrt{-1}\tau}$ 与 $q'=\ee^{2\pi\sqrt{-1}\tau'}$ 给出，其中 $\tau,\tau'$ 是两个规范因子的复化规范耦合。

来自每个不动点的贡献 $Z(\vec a,\vec Y;\vec b,\vec W)$，由各个多重态贡献的乘积给出。$\SU(N)$ 的矢量多重态贡献为
\begin{equation}
\resizebox{0.92\textwidth}{!}{$\displaystyle
Z^{\text{vector},\SU(N)}_{\text{instanton}}(\vec a,\vec Y)=\frac{1}{\prod_{i,j=1}^{N}\prod_{s\in Y_i}E(a_i-a_j,Y_i,Y_j,s)\prod_{t\in Y_j}\big(\epsilon_1+\epsilon_2-E(a_j-a_i,Y_j,Y_i,t)\big)}
$}
\tag{4.23}
\end{equation}
而 $\SU(M)$ 的贡献可以类似地得到。质量为 $m$ 的双基表示的贡献为
\begin{equation}
\begin{aligned}
&Z^{\text{hyper,bifundamental}}_{\text{instanton}}(\vec a,\vec Y;\vec b,\vec W;m)\\
&\quad=\prod_{i}^{N}\prod_{j}^{M}\Bigg[\prod_{s\in Y_i}\Big(E(a_i-b_j,Y_i,W_j,s)-m-\frac{\epsilon_1+\epsilon_2}{2}\Big)\\
&\qquad\qquad\times\prod_{t\in W_j}\Big(E(b_j-a_i,W_j,Y_i,t)-m+\frac{\epsilon_1+\epsilon_2}{2}\Big)\Bigg].
\end{aligned}
\tag{4.24}
\end{equation}

% [Tutor Note: p.27, (4.23)(4.24) 的读数方式：把 $\vec Y=(Y_1,\dots,Y_N)$ 里每个 $Y_i$ 看成一个「瞬子团簇」，
%  则 (4.23) 的分母是 $\prod_{i,j}$ 对所有配对 $(Y_i,Y_j)$ 的乘积，其中 $\prod_{s\in Y_i}E(a_i-a_j,Y_i,Y_j,s)$
%  里的盒子 $s$ 属于 $Y_i$、而臂长配 $Y_i$、leg长配 $Y_j$（见 (4.21) 的注释）；第二个因子
%  $\prod_{t\in Y_j}(\epsilon_1+\epsilon_2-E(a_j-a_i,Y_j,Y_i,t))$ 是把 $i,j$ 互换后的对应项。注意这里对 $j$ 的
%  求和是 $j=1,\dots,N$、没有限制 $j\neq i$——$i=j$ 的项正好给出 $E(a_i-a_i,Y_i,Y_i,s)$，不是零，不能漏掉。
%  (4.24) 的双基表示贡献同样是 $\prod_i\prod_j$；伴随表示是令 $\vec a=\vec b,\vec Y=\vec W$ 的特例。]

% [Tutor Note: p.27, (4.24) 中两个 $1/2$ 前的符号相反：第一组是 $-m-(\epsilon_1+\epsilon_2)/2$，
%  第二组是 $-m+(\epsilon_1+\epsilon_2)/2$。把它们抄成同号是本段最常见的错误，可用 p.28 的
%  (4.26) 反查：只有当第一组为 $-m-\cdots$、第二组为 $-m+\cdots$ 时，(4.25) 分子在
%  $m=(\epsilon_1-\epsilon_2)/2$ 处才恰好给出 $\Gamma_2((a_i-a_j)+\epsilon_1)\Gamma_2((a_i-a_j)+\epsilon_2)$。]

来自 $\SU(N)$ 的伴随超多重态的贡献可以通过令 $\vec a=\vec b$ 与 $\vec Y=\vec W$ 得到。类似地，来自 $\SU(N)$ 的 $N_f$ 个基础超多重态的贡献可以通过令 $M=N_f$、把 $\SU(M)$ 部分视作背景规范场并令 $W_i=0$ 得到。此时 $\vec b$ 成为质量参数中属于 $\SU(N_f)$ 的部分。作为例子，
% [原文 p.28]
具有质量为 $m$ 的伴随超多重态的 $\SU(N)$ 规范理论的配分函数为
\begin{equation}
\resizebox{0.9\textwidth}{!}{$\displaystyle
\begin{aligned}
Z(\tau,m)&=\frac{1}{N!}\int\dd^{N-1}a\,\Big|q^{-\frac{1}{\epsilon_1\epsilon_2}\sum a_i^2}\Big|^2\\
&\quad\times\prod_{i,j}\frac{\Gamma_2((a_i-a_j)+m+\frac{\epsilon_1+\epsilon_2}{2}|\epsilon_1,\epsilon_2)\Gamma_2((a_i-a_j)-m+\frac{\epsilon_1+\epsilon_2}{2}|\epsilon_1,\epsilon_2)}{\Gamma_2((a_i-a_j)+\epsilon_1+\epsilon_2|\epsilon_1,\epsilon_2)\Gamma_2((a_i-a_j)|\epsilon_1,\epsilon_2)}\\
&\quad\times\Bigg|\sum_{\vec Y}q^{|\vec Y|}\prod_{i,j}\frac{\prod_{s\in Y_i}\big(E(a_i-a_j,Y_i,Y_j,s)-m-\frac{\epsilon_1+\epsilon_2}{2}\big)\prod_{t\in Y_j}\big(E(a_j-a_i,Y_j,Y_i,t)-m+\frac{\epsilon_1+\epsilon_2}{2}\big)}{\prod_{s\in Y_i}E(a_i-a_j,Y_i,Y_j,s)\prod_{t\in Y_j}\big(\epsilon_1+\epsilon_2-E(a_j-a_i,Y_j,Y_i,t)\big)}\Bigg|^2
\end{aligned}
$}
\tag{4.25}
\end{equation}

% [Tutor Note: p.28, (4.25) 的四项因子与来源一一对应：整体 $|\cdot|^2$ 是 (4.15) 的
%  $Z(a)\overline{Z(a)}$（南北半球）；第一行 $|q^{-\frac{1}{\epsilon_1\epsilon_2}\sum a_i^2}|^2$
%  来自 (4.17)（$a$ 纯虚，故 $\sum a_i^2<0$）；第二行是 (4.18)(4.19) 的单圈因子之比
%  （分子两个 $\Gamma_2$ 来自伴随超多重态、分母两个来自矢量多重态）；第三行是 (4.23) 与 (4.24)
%  在 $\vec b=\vec a,\vec W=\vec Y$ 时的瞬子部分。抄写时最容易漏的是第三行分母里的
%  $\prod_{t\in Y_j}(\epsilon_1+\epsilon_2-E(\cdot))$——这一项来自 (4.23) 的第二个乘积，不是分子。]

这个表达式很可怕，但仍然相当显式。理论的内容物与规范群为 $\SU(N)$ 的 $\cN=4$ 超杨–米尔斯相同，而这个超引力背景上的超对称被认为会在 $m=\pm(\epsilon_1-\epsilon_2)/2$ 时增强。此时表达式大大简化 [29]。在这种情况下，瞬子修正是平凡的，因为当 $|Y|\ge1$ 时，存在一个盒子使分子中至少一个因子为零。

% [Tutor Note: p.28, 原文 "when $|Y|\ge1$" 处按 p.26 的定义应写作 $|\vec Y|\ge1$（瞬子数 = 各图盒子数之和）；
%  已用 700 dpi 裁图复核，原文确实只写了 $|Y|$、没有箭头，属排印遗漏，正文照原样排印。
%  另："Hermitean" 按通行拼法应为 "Hermitian"（同样 700 dpi 复核）。
%  提示："the matter content of the the-ory" 不是笔误——那是 "the theory" 在行末的连字断行
%  （"the-" + 换行的 "ory"），不要当成重复的 "the"。]

% [Tutor Note: p.28, 原文「存在一个盒子使至少一个分子因子为零」的具体机制（补出）：
%  下面以 $m=-\frac{\epsilon_1-\epsilon_2}{2}$ 为例核验。对每张非空的 Young 图 $Y_i$，令
%  $w_i=\lambda'_1$ 为其宽度，取它最右一列最上面的那个盒子 $s=(w_i,\lambda_{w_i})$。由列高不增知
%  所有列的列高都 $\ge\lambda_{w_i}$，故 $L_{Y_i}(s)=\lambda'_{\lambda_{w_i}}-w_i=w_i-w_i=0$，
%  同时 $A_{Y_i}(s)=\lambda_{w_i}-\lambda_{w_i}=0$；于是在对角配对 $(i,i)$ 上
%  $E(a_i-a_i,Y_i,Y_i,s)=-\epsilon_1\cdot0+\epsilon_2(0+1)=\epsilon_2$，而 (4.25) 分子第一个因子
%  $\epsilon_2-m-\frac{\epsilon_1+\epsilon_2}{2}=\epsilon_2-\frac{\epsilon_2-\epsilon_1}{2}-\frac{\epsilon_1+\epsilon_2}{2}=0$。
%  故凡 $|\vec Y|\ge1$ 的项全为零，只剩 $|\vec Y|=0$ 的项（即 1），瞬子修正因此是平凡的。
%  注意这里 $s$ 落在 $Y_i$ 内部，所以 $A,L$ 都是非负整数——「$s$ 可以在图外」这一约定是给
%  $i\neq j$ 的配对用的，本论证用不到。]

对每个 $(i,j)$ 的选择，单圈因子给出
\begin{equation}
\frac{\Gamma_2((a_i-a_j)+\epsilon_1|\epsilon_1,\epsilon_2)\Gamma_2((a_i-a_j)+\epsilon_2|\epsilon_1,\epsilon_2)}{\Gamma_2((a_i-a_j)|\epsilon_1,\epsilon_2)\Gamma_2((a_i-a_j)+\epsilon_1+\epsilon_2|\epsilon_1,\epsilon_2)}=a_i-a_j.
\tag{4.26}
\end{equation}

% [Tutor Note: p.28, (4.26) 是 Barnes 双 Gamma 函数的功能方程
%  $\Gamma_2(x+\epsilon_1)\Gamma_2(x+\epsilon_2)/[\Gamma_2(x)\Gamma_2(x+\epsilon_1+\epsilon_2)]=x$
%  在 $x=a_i-a_j$ 处的写法。把 $m=\pm(\epsilon_1-\epsilon_2)/2$ 代进 (4.25) 第二行即可看出它正是
%  (4.18)(4.19) 的单圈因子在该质量处的值：取 $m=(\epsilon_1-\epsilon_2)/2$，则分子第一个
%  $\Gamma_2$ 的自变量是 $(a_i-a_j)+\frac{\epsilon_1-\epsilon_2}{2}+\frac{\epsilon_1+\epsilon_2}{2}=(a_i-a_j)+\epsilon_1$，
%  第二个是 $(a_i-a_j)-\frac{\epsilon_1-\epsilon_2}{2}+\frac{\epsilon_1+\epsilon_2}{2}=(a_i-a_j)+\epsilon_2$，
%  分母恰为 $(a_i-a_j)+\epsilon_1+\epsilon_2$ 与 $(a_i-a_j)$——与 (4.26) 完全一致。]

于是积分归结为
\begin{equation}
\resizebox{0.9\textwidth}{!}{$\displaystyle
Z\Big(\tau,\frac{\epsilon_1-\epsilon_2}{2}\Big)=\frac{1}{N!}\int\dd^{N-1}a\prod_{i,j}(a_i-a_j)\Big|q^{-\frac{1}{\epsilon_1\epsilon_2}\sum a_i^2}\Big|^2=\int\dd^{N^2-1}A\,|q|^{-\frac{2}{\epsilon_1\epsilon_2}\tr AA^\dagger}
$}
\tag{4.27}
\end{equation}
也就是说，这正是一个厄米随机矩阵模型。

% [Tutor Note: p.28, 从 (4.25) 到 (4.27) 走了两步：(i) $m=\pm(\epsilon_1-\epsilon_2)/2$ 时瞬子求和只剩
%  $|\vec Y|=0$ 的项（见上一条批注的机制），于是整个瞬子分数化为 1；(ii) 单圈因子由 (4.26) 化为
%  $a_i-a_j$，与测度里本就有的 $\prod_{i,j}$ 相乘，得到 Vandermonde 因子的平方：
%  $\prod_{i\ne j}(a_i-a_j)=\prod_{i<j}(a_i-a_j)(a_j-a_i)=(-1)^{N(N-1)/2}\prod_{i<j}(a_i-a_j)^2$
%  （整体符号可被测度或 $a\to-a$ 的换元吸收）。最后一步把对角化的 $\vec a$ 换成 $\U(N)$ 的厄米矩阵 $A$：
%  $d^{N^2-1}A$ 是去掉迹部分（即 $\U(1)$ 部分）后的测度，$A$ 的矩阵元给出高斯权重，
%  故这是标准的厄米矩阵模型——本综述把它与 Liouville 理论联系起来（见后文 §5.2）。]

% [Tutor Note: p.28, (4.27) 中的 $\prod_{i,j}(a_i-a_j)$ 若按 $i,j=1,\dots,N$ 直读，$i=j$ 时因子为 0，
%  整个积分会为零。按上一条批注的 Vandermonde 结构，这里应是 $\prod_{i<j}$（或 $\prod_{i\ne j}$），
%  即 $\pm\prod_{i<j}(a_i-a_j)^2$。已用 900 dpi 裁图复核：原文确实只印了下标 $i,j$、没有写 $i<j$，
%  属原文的记号疏漏（同 (4.12)）；正文照原样排印，计算时按 $\prod_{i\ne j}$ 理解。]

这一点最初是在 [30,31] 中观察到的，而 [6] 的 $S^4$ 局域化最初正是为了从量子场论上导出这一事实而提出的。
